\maketitle

\section{Keller Groupoids and Fiber Configuration
Geometry}\label{keller-groupoids-and-fiber-configuration-geometry}

\textbf{Working draft / research-state document}\\
\textbf{Snapshot:} 2026-09-14

This document reframes the existing Keller-threefold work around the
\textbf{Keller groupoid} attached to an arbitrary Keller map. The goal
is to extract everything that is formal or general from the one known
nontrivial example before specializing to the marked-cubic degree-three
map.

The intended division is:

\begin{enumerate}
\def\labelenumi{\arabic{enumi}.}
\tightlist
\item
  general machinery valid for every complex Keller map;
\item
  finite-degree and monodromy consequences;
\item
  topology of the Čech nerve;
\item
  finite-stabilization principles;
\item
  the marked-cubic example as the first worked nontrivial Keller
  groupoid;
\item
  a research program aimed particularly at dimension two.
\end{enumerate}

The document distinguishes \textbf{proved here}, \textbf{standard
imported machinery}, and \textbf{open program}. Nothing in the
open-program sections is being promoted to a theorem.

\begin{center}\rule{0.5\linewidth}{0.5pt}\end{center}

\section{1. Keller maps and the relation
groupoid}\label{keller-maps-and-the-relation-groupoid}

\subsection{Definition 1.1 --- Keller
map}\label{definition-1.1-keller-map}

A \textbf{complex Keller map} is a polynomial map

\[
f\colon X=\mathbb A^n_{\mathbb C}\longrightarrow
Y=\mathbb A^n_{\mathbb C}
\]

whose Jacobian determinant is a nonzero constant.

Equivalently, \(f\) is étale. In particular it is flat, unramified,
locally quasi-finite, separated, and open.

Write

\[
U=f(X)\subseteq Y.
\]

Since an étale map is open, \(U\) is an open subvariety and

\[
f\colon X\longrightarrow U
\]

is a surjective étale morphism.

For a dominant Keller map, let

\[
d=[\mathbb C(X):\mathbb C(Y)]
\]

be its \textbf{generic degree}.

\begin{center}\rule{0.5\linewidth}{0.5pt}\end{center}

\subsection{Definition 1.2 --- Keller relation
groupoid}\label{definition-1.2-keller-relation-groupoid}

Define

\[
R_f=X\times_Y X
=\{(x_0,x_1):f(x_0)=f(x_1)\}.
\]

The two projections give source and target maps

\[
s,t\colon R_f\rightrightarrows X,
\qquad
s(x_0,x_1)=x_0,\quad t(x_0,x_1)=x_1.
\]

Together with

\[
e(x)=(x,x),\qquad
i(x_0,x_1)=(x_1,x_0),
\]

and

\[
(x_0,x_1)\circ(x_1,x_2)=(x_0,x_2),
\]

this is a \textbf{thin étale groupoid}: there is at most one arrow
between any ordered pair of objects, and an arrow exists exactly when
the two points have the same \(f\)-image.

We call this the \textbf{Keller groupoid} of \(f\), denoted
\(\mathcal G_f\).

Because both maps \(X\to Y\) factor through \(U\),

\[
X\times_Y X=X\times_U X.
\]

Thus \(\mathcal G_f\) is exactly the equivalence-relation groupoid of
the surjective étale map \(X\to U\).

\begin{center}\rule{0.5\linewidth}{0.5pt}\end{center}

\subsection{Proposition 1.3 --- The diagonal is a connected
component}\label{proposition-1.3-the-diagonal-is-a-connected-component}

The identity section

\[
\Delta_X\colon X\hookrightarrow R_f
\]

is both open and closed.

\subsubsection{Proof}\label{proof}

A Keller map is unramified, hence its diagonal is an open immersion. It
is also separated, hence its diagonal is a closed immersion. Therefore
the diagonal is open and closed.

This is the standard unramified/separated diagonal theorem; see Stacks
Project Tags \texttt{02GE} and \texttt{024T}.

Consequently,

\[
R_f=\Delta_X\sqcup \operatorname{Conf}_2(f),
\]

where

\[
\operatorname{Conf}_2(f)=\{(x_0,x_1):f(x_0)=f(x_1),\ x_0\ne x_1\}.
\]

The space \(\operatorname{Conf}_2(f)\) is the first \textbf{fiber
configuration variety}.

\begin{center}\rule{0.5\linewidth}{0.5pt}\end{center}

\subsection{Corollary 1.4 --- Automorphisms are exactly the
collision-free
case}\label{corollary-1.4-automorphisms-are-exactly-the-collision-free-case}

For a complex Keller map,

\[
\operatorname{Conf}_2(f)=\varnothing
\quad\Longleftrightarrow\quad
f\text{ is injective}
\quad\Longleftrightarrow\quad
f\in\operatorname{Aut}(\mathbb A^n).
\]

Hence the Jacobian conjecture in dimension \(n\) can be restated as:

\begin{quote}
Every Keller groupoid on \(\mathbb A^n\) is the unit groupoid.
\end{quote}

The last implication uses the standard characteristic-zero
injective-polynomial-map theorem (Ax--Grothendieck).

This is the most basic groupoid reformulation: \textbf{a counterexample
is exactly a Keller groupoid with nonempty fiber configuration space.}

\begin{center}\rule{0.5\linewidth}{0.5pt}\end{center}

\section{2. The Čech nerve}\label{the-ux10dech-nerve}

\subsection{Definition 2.1 --- Higher Keller fiber
products}\label{definition-2.1-higher-keller-fiber-products}

For \(r\ge1\), put

\[
K_r(f)
=
\underbrace{X\times_Y X\times_Y\cdots\times_Y X}_{r\text{ factors}}.
\]

Thus

\[
K_r(f)
=
\{(x_1,\dots,x_r):f(x_1)=\cdots=f(x_r)\}.
\]

We index this so that

\[
K_1(f)=X,\qquad K_2(f)=R_f.
\]

Deleting a coordinate gives a face map; repeating a coordinate gives a
degeneracy map. These spaces form the Čech nerve of \(X\to U\), or
equivalently the nerve of the Keller groupoid.

\begin{center}\rule{0.5\linewidth}{0.5pt}\end{center}

\subsection{\texorpdfstring{Proposition 2.2 --- Every level is a smooth
affine
\(n\)-fold}{Proposition 2.2 --- Every level is a smooth affine n-fold}}\label{proposition-2.2-every-level-is-a-smooth-affine-n-fold}

For every \(r\ge1\):

\begin{enumerate}
\def\labelenumi{\arabic{enumi}.}
\tightlist
\item
  \(K_r(f)\) is affine and of finite type;
\item
  every projection \[
  p_i\colon K_r(f)\to X
  \] is étale;
\item
  \(K_r(f)\) is smooth of pure dimension \(n\).
\end{enumerate}

\subsubsection{Proof}\label{proof-1}

Fiber products of affine schemes over an affine scheme are affine. Each
projection is obtained by repeated base change from the étale map \(f\),
hence is étale. An étale map to the smooth \(n\)-fold \(X\) has smooth
\(n\)-dimensional source.

Thus the simplicial direction grows indefinitely while the algebraic
dimension never grows.

\begin{center}\rule{0.5\linewidth}{0.5pt}\end{center}

\subsection{Proposition 2.3 --- Canonical algebraic
parallelism}\label{proposition-2.3-canonical-algebraic-parallelism}

Every \(K_r(f)\), and every open-and-closed subvariety of it, has
trivial algebraic tangent bundle.

In particular, every fiber configuration variety introduced below is
algebraically parallelizable:

\[
T_{K_r(f)}\simeq\mathcal O_{K_r(f)}^{\oplus n}.
\]

\subsubsection{Proof}\label{proof-2}

For any projection \(p_i\colon K_r(f)\to X\), étaleness gives

\[
T_{K_r(f)}\simeq p_i^*T_X.
\]

Since \(T_{\mathbb A^n}\) is trivial, so is \(T_{K_r(f)}\).

This is a useful additional restriction on every variety appearing in a
Keller groupoid: they are not merely smooth affine \(n\)-folds, but
smooth affine \textbf{parallelizable} \(n\)-folds.

\begin{center}\rule{0.5\linewidth}{0.5pt}\end{center}

\subsection{Proposition 2.4 --- The infinite nerve is categorically
finite
data}\label{proposition-2.4-the-infinite-nerve-is-categorically-finite-data}

The simplicial object \(K_\bullet(f)\) is the nerve of a groupoid and
hence is \(2\)-coskeletal in the standard categorical sense.

Informally:

\begin{quote}
once \(K_1\), \(K_2\), and the source/target/composition/unit/inverse
maps are known, every higher \(K_r\) and every simplicial structure map
is determined by iterated fiber products.
\end{quote}

This is a \textbf{categorical stabilization} valid for every Keller map.

It must not be confused with the much stronger statement that the higher
\(K_r\) have no new isomorphism types of connected components. The
latter is a geometric question addressed below.

\begin{center}\rule{0.5\linewidth}{0.5pt}\end{center}

\section{3. Collision varieties and partition
decomposition}\label{collision-varieties-and-partition-decomposition}

\subsection{\texorpdfstring{Definition 3.1 --- Ordered \(k\)-fiber
configuration
variety}{Definition 3.1 --- Ordered k-fiber configuration variety}}\label{definition-3.1-ordered-k-fiber-configuration-variety}

For \(k\ge1\), define

\[
\operatorname{Conf}_k(f)=
\{(x_1,\dots,x_k)\in K_k(f):x_i\ne x_j\text{ for }i\ne j\}.
\]

Thus

\[
\operatorname{Conf}_1(f)=X,\qquad \operatorname{Conf}_2(f)=\text{ordered two-point configurations}.
\]

Because the conditions \(x_i=x_j\) are open-and-closed conditions in the
relevant fiber products, \(\operatorname{Conf}_k(f)\) is itself open and
closed in \(K_k(f)\).

Hence every \(\operatorname{Conf}_k(f)\) is a smooth affine
parallelizable \(n\)-fold, possibly disconnected or empty.

\begin{center}\rule{0.5\linewidth}{0.5pt}\end{center}

\subsection{Theorem 3.2 --- Partition
decomposition}\label{theorem-3.2-partition-decomposition}

Let \(\Pi(r)\) denote the set of set partitions of \(\{1,\dots,r\}\).
For \(\pi\in\Pi(r)\), let \(|\pi|\) be its number of blocks.

Then

\[
\boxed{
K_r(f)\simeq
\coprod_{\pi\in\Pi(r)} \operatorname{Conf}_{|\pi|}(f).
}
\]

More precisely, the summand corresponding to \(\pi\) is the locus on
which

\[
x_i=x_j
\quad\Longleftrightarrow\quad
i,j\text{ lie in the same block of }\pi.
\]

\subsubsection{Proof}\label{proof-3}

For each pair \(i,j\), the equalizer condition \(x_i=x_j\) is the
pullback of the diagonal of the unramified separated map \(f\), and is
therefore open and closed. Intersecting the required equalities and
inequalities for a fixed set partition produces an open-and-closed
stratum.

If \(\pi\) has \(k\) blocks, order those blocks by their least elements.
Reading the common coordinate value on each block gives an isomorphism
with \(\operatorname{Conf}_k(f)\).

\begin{center}\rule{0.5\linewidth}{0.5pt}\end{center}

\subsection{Corollary 3.3 --- Stirling-number
form}\label{corollary-3.3-stirling-number-form}

Let \(S(r,k)\) be the Stirling number of the second kind. Then

\[
\boxed{
K_r(f)\simeq
\coprod_{k\ge1} S(r,k)\,\operatorname{Conf}_k(f),
}
\]

where \(mZ\) means the disjoint union of \(m\) copies of \(Z\).

This partition decomposition is completely general. What is special in
the marked-cubic example is that several different
\(\operatorname{Conf}_k\)'s turn out to be isomorphic.

\begin{center}\rule{0.5\linewidth}{0.5pt}\end{center}

\section{4. Generic degree bounds the configuration
depth}\label{generic-degree-bounds-the-configuration-depth}

The next point is crucial: the generic degree gives a \emph{uniform
maximum} on fiber cardinality for an étale quasi-finite map.

\subsection{Lemma 4.1 --- No fiber has more points than the generic
degree}\label{lemma-4.1-no-fiber-has-more-points-than-the-generic-degree}

If \(f\) has generic degree \(d\), then

\[
\#f^{-1}(y)\le d
\]

for every \(y\in Y\).

\subsubsection{Proof sketch}\label{proof-sketch}

Suppose \(y\) has \(m\) distinct preimages \(x_1,\dots,x_m\). Because
\(f\) is étale and quasi-finite, after shrinking around \(y\) one can
choose pairwise disjoint neighborhoods of the \(x_i\) that persist as
\(m\) separate étale local sheets over a common target neighborhood.
Hence a generic point of that neighborhood has at least \(m\) preimages.
By definition the generic number is \(d\), so \(m\le d\).

\begin{center}\rule{0.5\linewidth}{0.5pt}\end{center}

\subsection{Corollary 4.2 --- Finite configuration
palette}\label{corollary-4.2-finite-configuration-palette}

\[
\operatorname{Conf}_k(f)=\varnothing
\qquad(k>d).
\]

Therefore

\[
\boxed{
K_r(f)\simeq
\coprod_{k=1}^{\min(r,d)}
S(r,k)\,\operatorname{Conf}_k(f).
}
\]

In particular:

\begin{quote}
\textbf{Every finite-degree Keller groupoid has only finitely many
geometric fiber configuration spaces
\(\operatorname{Conf}_1,\dots,\operatorname{Conf}_d\), even though its
Čech nerve has infinitely many levels.}
\end{quote}

This is stronger than the conditional stabilization principle we
initially expected: all connected-component isomorphism types occurring
in the entire Keller groupoid necessarily occur by level \(d\).

The special example below stabilizes \emph{earlier} than this universal
bound.

\begin{center}\rule{0.5\linewidth}{0.5pt}\end{center}

\section{5. The fiber-cardinality
filtration}\label{the-fiber-cardinality-filtration}

\subsection{Definition 5.1}\label{definition-5.1}

For \(1\le k\le d\), define

\[
U_{\ge k}
=
\{y\in Y:\#f^{-1}(y)\ge k\}.
\]

Also put

\[
U_m
=
\{y\in Y:\#f^{-1}(y)=m\}
\qquad(0\le m\le d).
\]

Thus \(U_{\ge1}=U=f(X)\), and \(U_0=Y\setminus U\).

\begin{center}\rule{0.5\linewidth}{0.5pt}\end{center}

\subsection{Proposition 5.2 --- Collision images recover the
filtration}\label{proposition-5.2-collision-images-recover-the-filtration}

The common-image map

\[
q_k\colon \operatorname{Conf}_k(f)\to Y,
\qquad
(x_1,\dots,x_k)\mapsto f(x_1),
\]

is étale and has image

\[
q_k(\operatorname{Conf}_k(f))=U_{\ge k}.
\]

Consequently \(U_{\ge k}\) is open and

\[
U_{\ge1}\supseteq U_{\ge2}\supseteq\cdots\supseteq U_{\ge d}.
\]

Over a point \(y\in U_m\), the fiber of \(q_k\) has cardinality

\[
(m)_k
=
m(m-1)\cdots(m-k+1)
=
\frac{m!}{(m-k)!}.
\]

\subsubsection{Proof}\label{proof-4}

The map \(q_k=f\circ p_1\) is a composition of étale maps, hence étale
and therefore open. Its geometric fiber is exactly the set of ordered
\(k\)-tuples of distinct elements of the finite set \(f^{-1}(y)\).

\begin{center}\rule{0.5\linewidth}{0.5pt}\end{center}

\subsection{Proposition 5.3 --- The full-fiber locus is the proper
locus}\label{proposition-5.3-the-full-fiber-locus-is-the-proper-locus}

The locus \(U_{\ge d}=U_d\) is exactly the locus over which \(f\) is
proper. Equivalently, if \(S_f\) denotes the nonproperness set,

\[
\boxed{
S_f=Y\setminus U_d.
}
\]

\subsubsection{Proof sketch}\label{proof-sketch-1}

If \(f\) is proper in a neighborhood of \(y\), then it is finite étale
there. Its degree is locally constant and equals the generic degree
\(d\), so the fiber has \(d\) points.

Conversely, suppose the fiber over \(y\) has \(d\) points. Étaleness
gives \(d\) disjoint local sheets over a sufficiently small common
target neighborhood. By Lemma 4.1 there can be no additional sheets
anywhere over that neighborhood. Hence the full inverse image of that
neighborhood is the union of those \(d\) finite étale sheets, so \(f\)
is finite, hence proper, there.

Thus the groupoid recovers not just the image but the entire fiber-drop
filtration and the nonproperness locus.

\begin{center}\rule{0.5\linewidth}{0.5pt}\end{center}

\section{6. The configuration tower}\label{the-configuration-tower}

Deleting the last coordinate gives

\[
\rho_k\colon \operatorname{Conf}_k(f)\longrightarrow C_{k-1}(f).
\]

\subsection{Proposition 6.1}\label{proposition-6.1}

The map \(\rho_k\) is étale. If a point of \(C_{k-1}\) lies over an
\(m\)-point fiber of \(f\), then its \(\rho_k\)-fiber contains exactly

\[
m-k+1
\]

points.

Generically,

\[
\deg(\rho_k)=d-k+1.
\]

\begin{center}\rule{0.5\linewidth}{0.5pt}\end{center}

\subsection{Proposition 6.2 --- Top collision open
immersion}\label{proposition-6.2-top-collision-open-immersion}

The top forgetful map

\[
\rho_d\colon \operatorname{Conf}_d(f)\to \operatorname{Conf}_{d-1}(f)
\]

is an open immersion.

Its image is exactly the locus of ordered \((d-1)\)-tuples lying over a
\(d\)-point fiber.

Therefore:

\[
\boxed{
\operatorname{Conf}_d(f)\simeq \operatorname{Conf}_{d-1}(f)
\quad\Longleftrightarrow\quad
f\text{ has no fiber of cardinality }d-1.
}
\]

\subsubsection{Proof}\label{proof-5}

The map is étale. A \((d-1)\)-tuple can have at most one additional
point, because no fiber contains more than \(d\) points. Thus \(\rho_d\)
is universally injective and étale, hence an open immersion.

It is surjective exactly when every fiber supporting a \((d-1)\)-tuple
has a remaining \(d\)-th point, i.e.~exactly when there is no
\((d-1)\)-point fiber.

This is the abstract mechanism behind the marked-cubic identity
\(\operatorname{Conf}_3\simeq\operatorname{Conf}_2\).

\begin{center}\rule{0.5\linewidth}{0.5pt}\end{center}

\section{7. Monodromy and connected
components}\label{monodromy-and-connected-components}

Let

\[
V=U_d=Y\setminus S_f.
\]

Then

\[
f^{-1}(V)\to V
\]

is a finite étale cover of degree \(d\).

Choose a geometric base point and write

\[
G\le S_d
\]

for its monodromy group. Since \(f^{-1}(V)\) is a nonempty open subset
of the irreducible variety \(X\), it is irreducible, hence the action of
\(G\) on the \(d\) sheets is transitive.

\begin{center}\rule{0.5\linewidth}{0.5pt}\end{center}

\subsection{Theorem 7.1 --- Collision components are monodromy
orbits}\label{theorem-7.1-collision-components-are-monodromy-orbits}

Over \(V\),

\[
\operatorname{Conf}_k(f)|_V\to V
\]

is the finite étale cover associated with the \(G\)-set

\[
\operatorname{Inj}(\{1,\dots,k\},\{1,\dots,d\}).
\]

Consequently, connected components of \(\operatorname{Conf}_k(f)\) are
in bijection with \(G\)-orbits on ordered \(k\)-tuples of distinct
sheets.

In particular:

\[
\operatorname{Conf}_k(f)\text{ is irreducible}
\quad\Longleftrightarrow\quad
G\text{ is }k\text{-transitive}.
\]

\subsubsection{Reason}\label{reason}

Every global component meets the dense proper locus \(V\); on a smooth
variety connected components are irreducible components, and an open
subset of an irreducible component remains irreducible. Thus the generic
finite-cover components and the global components correspond.

\begin{center}\rule{0.5\linewidth}{0.5pt}\end{center}

\subsection{Corollary 7.2 --- The top configuration is the Galois
closure}\label{corollary-7.2-the-top-configuration-is-the-galois-closure}

The fiber of \(\operatorname{Conf}_d|_V\) is the set of all orderings of
the \(d\) sheets, hence has \(d!\) elements.

The action of \(G\) on these orderings is free. Therefore:

\begin{itemize}
\tightlist
\item
  \(\operatorname{Conf}_d|_V\) has \[
  \frac{d!}{|G|}
  \] connected components;
\item
  each component is a finite étale Galois cover of \(V\) with deck group
  \(G\);
\item
  each component is a geometric realization of the Galois closure of the
  function-field extension.
\end{itemize}

Thus the top fiber configuration space makes the monodromy/Galois
closure intrinsic to the Keller groupoid.

\begin{center}\rule{0.5\linewidth}{0.5pt}\end{center}

\subsection{Corollary 7.3 --- Normal extensions are groupoid-simple and
forbidden for
counterexamples}\label{corollary-7.3-normal-extensions-are-groupoid-simple-and-forbidden-for-counterexamples}

If the original function-field extension is Galois, then the monodromy
action on the \(d\) sheets is regular, so \(|G|=d\).

Campbell's classical normal-extension criterion says that a complex
Keller map with normal function-field extension is invertible. Therefore
a nontrivial Keller groupoid must have

\[
|G|>d.
\]

For \(d=2\), every transitive degree-two group is regular. Hence there
is no degree-two complex Keller counterexample.

For \(d=3\), the first non-Galois transitive possibility is \(S_3\);
this is exactly the monodromy of the marked-cubic example.

This explains structurally why generic degree three is the first
possible degree for the known construction.

\begin{center}\rule{0.5\linewidth}{0.5pt}\end{center}

\section{8. Quotients of fiber configuration
spaces}\label{quotients-of-fiber-configuration-spaces}

The symmetric group \(S_k\) acts freely on \(\operatorname{Conf}_k(f)\)
by permuting coordinates.

The affine quotient

\[
B_k(f)=\operatorname{Conf}_k(f)/S_k
\]

parametrizes unordered \(k\)-element subsets of a fiber.

Over \(U_m\), its fiber has cardinality

\[
\binom{m}{k}.
\]

At the top level there is a unique \(d\)-element subset of a
\(d\)-element fiber, so

\[
\boxed{
B_d(f)\simeq U_d=Y\setminus S_f.
}
\]

Thus the proper locus itself is the symmetric quotient of the top
configuration variety.

Over the proper locus there is also the usual complement duality

\[
B_k|_V\simeq B_{d-k}|_V,
\]

sending a subset of the \(d\) sheets to its complement.

\begin{center}\rule{0.5\linewidth}{0.5pt}\end{center}

\section{9. Additive invariants and Stirling
calculus}\label{additive-invariants-and-stirling-calculus}

The partition decomposition immediately gives identities in the
Grothendieck ring of varieties:

\[
\boxed{
[K_r(f)]
=
\sum_{k=1}^d S(r,k)[\operatorname{Conf}_k(f)].
}
\]

No monodromy assumption is involved.

The same formula holds for every additive invariant \(I\) of complex
algebraic varieties:

\[
\boxed{
I(K_r)
=
\sum_{k=1}^d S(r,k)I(\operatorname{Conf}_k).
}
\]

Examples include Euler characteristic and any invariant obtained from a
Grothendieck-ring realization for which the required additivity is
available.

\begin{center}\rule{0.5\linewidth}{0.5pt}\end{center}

\subsection{Corollary 9.1 --- Exponential generating
function}\label{corollary-9.1-exponential-generating-function}

Formally,

\[
\sum_{r\ge1}I(K_r)\frac{t^r}{r!}
=
\sum_{k=1}^d
I(\operatorname{Conf}_k)\frac{(e^t-1)^k}{k!}.
\]

Thus the entire infinite sequence of levelwise additive invariants is
determined by the first \(d\) configuration invariants.

\begin{center}\rule{0.5\linewidth}{0.5pt}\end{center}

\subsection{Corollary 9.2 --- Stirling
inversion}\label{corollary-9.2-stirling-inversion}

The configuration invariants can conversely be recovered from the level
invariants:

\[
I(\operatorname{Conf}_k)
=
\sum_{r=1}^k s(k,r)I(K_r),
\]

where \(s(k,r)\) are the signed Stirling numbers of the first kind.

So \(K_\bullet\) and \(C_\bullet\) carry the same additive information,
expressed in two natural bases.

\begin{center}\rule{0.5\linewidth}{0.5pt}\end{center}

\subsection{Corollary 9.3 --- Linear
recurrence}\label{corollary-9.3-linear-recurrence}

Since \(S(r,k)\) is a linear combination of

\[
1^r,2^r,\dots,k^r,
\]

the scalar sequence

\[
a_r=I(K_r)
\]

is a linear combination of

\[
1^r,2^r,\dots,d^r.
\]

Hence it satisfies the constant-coefficient recurrence whose
characteristic polynomial is

\[
\boxed{
\prod_{j=1}^d(T-j).
}
\]

Equivalently, its ordinary generating function is rational with
denominator dividing

\[
\prod_{j=1}^d(1-jz).
\]

This is one concrete sense in which the infinite Keller nerve is
automatically finite-state at the level of additive invariants.

\begin{center}\rule{0.5\linewidth}{0.5pt}\end{center}

\section{10. Finite appearance of connected-component
types}\label{finite-appearance-of-connected-component-types}

This section records the stabilization principle that motivated the
present reframing.

Let

\[
\operatorname{Conf}_k=\coprod_j a_{k,j}Z_j
\]

where the \(Z_j\)'s are representatives of connected-component
isomorphism types and \(a_{k,j}\) their multiplicities.

Because \(k\le d\), there are only finitely many \(Z_j\)'s.

\subsection{Theorem 10.1 --- Finite geometric
palette}\label{theorem-10.1-finite-geometric-palette}

Every connected component of every \(K_r(f)\) is isomorphic to a
connected component of one of

\[
\operatorname{Conf}_1(f),\dots,\operatorname{Conf}_d(f).
\]

Therefore all connected-component isomorphism types appear by finite
step \(d\).

More generally, suppose all component types happen already to appear by
some step \(m<d\). Then every later level is a disjoint union of that
finite palette.

If \(N_j(r)\) is the number of copies of \(Z_j\) in \(K_r\), then

\[
\boxed{
N_j(r)=\sum_{k=1}^d a_{k,j}S(r,k).
}
\]

Consequently every multiplicity sequence \(N_j(r)\):

\begin{itemize}
\tightlist
\item
  is a linear combination of \(1^r,\dots,d^r\);
\item
  satisfies the recurrence \(\prod_{i=1}^d(T-i)\);
\item
  has a rational ordinary generating function.
\end{itemize}

So an early geometric collapse has immediate, testable combinatorial
consequences at \textbf{all} later levels.

\begin{center}\rule{0.5\linewidth}{0.5pt}\end{center}

\subsection{Definition 10.2 --- Geometric generation
depth}\label{definition-10.2-geometric-generation-depth}

Define the \textbf{geometric generation depth}

\[
g(f)
\]

to be the least \(m\) such that every connected-component isomorphism
type appearing in every \(\operatorname{Conf}_k\) already appears among
\(\operatorname{Conf}_1,\dots,\operatorname{Conf}_m\).

Always

\[
1\le g(f)\le d.
\]

For an automorphism, \(g(f)=1\).

For the marked-cubic example, \(g(F)=2\).

This is a first candidate for a quantitative notion of
\textbf{configuration complexity}.

\begin{center}\rule{0.5\linewidth}{0.5pt}\end{center}

\subsection{Stronger notion: finite morphism
type}\label{stronger-notion-finite-morphism-type}

Object-type stabilization does not by itself determine all topology of
the realization, because the face maps matter.

A stronger condition is:

\begin{quote}
up to chosen component isomorphisms, only finitely many morphism types
occur among restrictions of face and degeneracy maps.
\end{quote}

Call this \textbf{finite morphism type}.

Under finite morphism type, applying any functor such as homology turns
the entire infinite simplicial object into a finite-state system:
finitely many modules and finitely many linear maps occur, with
combinatorially varying multiplicities.

The marked-cubic groupoid satisfies an exceptionally strong version of
this condition; see Section 14.

This suggests a hierarchy:

\[
\text{finite generic degree}
\Rightarrow
\text{finite object palette}
\]

always, while

\[
\text{early palette stabilization}
\]

and

\[
\text{finite morphism type}
\]

measure extra simplicity.

\begin{center}\rule{0.5\linewidth}{0.5pt}\end{center}

\section{11. Euler characteristics of fiber-cardinality
strata}\label{euler-characteristics-of-fiber-cardinality-strata}

Let

\[
a_m=\chi(U_m),
\qquad
b_k=\chi(\operatorname{Conf}_k).
\]

Since

\[
\operatorname{Conf}_k|_{U_m}\to U_m
\]

is a finite topological covering of degree \((m)_k\),

\[
\boxed{
b_k
=
\sum_{m=k}^d (m)_k\,a_m.
}
\]

This is a triangular system with diagonal entries \(k!\). Therefore the
Euler characteristics of the exact fiber-cardinality strata are
recoverable from the collision Euler characteristics.

Since

\[
\chi(Y)=\chi(\mathbb A^n)=1,
\]

we also have

\[
1
=
a_0+\sum_{m=1}^d a_m.
\]

And because

\[
b_1=\chi(X)=1
=
\sum_{m=1}^d m\,a_m,
\]

subtracting gives

\[
\boxed{
a_0
=
\sum_{m=1}^d (m-1)a_m.
}
\]

Thus the Euler characteristic of the omitted locus is encoded by the
fiber-drop strata.

This is potentially useful in dimension two, where the omitted locus of
a Keller map is finite.

\begin{center}\rule{0.5\linewidth}{0.5pt}\end{center}

\section{12. Topological realization of the Keller
groupoid}\label{topological-realization-of-the-keller-groupoid}

The following is standard descent/topological-groupoid machinery rather
than a new theorem.

Because

\[
f\colon X^{an}\to U^{an}
\]

is a surjective local homeomorphism, the topological Čech nerve

\[
K_\bullet(f)^{an}
\]

is a simplicial resolution of \(U^{an}\).

Its geometric realization satisfies, in the standard homotopy-descent
sense,

\[
\boxed{
|K_\bullet(f)^{an}|\simeq U^{an}.
}
\]

Equivalently, the quotient of the étale equivalence-relation groupoid
recovers the image \(U\).

There is a homological descent spectral sequence

\[
\boxed{
E^1_{p,q}
=
H_q(K_{p+1}(f)^{an};\mathbb Z)
\Longrightarrow
H_{p+q}(U^{an};\mathbb Z),
}
\]

with \(d^1\) the alternating sum of the face maps.

\begin{center}\rule{0.5\linewidth}{0.5pt}\end{center}

\subsection{Proposition 12.1 --- Vertical homological
boundedness}\label{proposition-12.1-vertical-homological-boundedness}

Every \(K_r(f)\) is a smooth affine complex \(n\)-fold. By the
Andreotti--Frankel theorem it has the homotopy type of a real CW complex
of dimension at most \(n\).

Therefore

\[
H_q(K_r;\mathbb Z)=0
\qquad(q>n).
\]

The descent spectral sequence is thus an infinite horizontal strip of
only \(n+1\) rows.

This is another general finite-complexity phenomenon.

\begin{center}\rule{0.5\linewidth}{0.5pt}\end{center}

\subsection{Corollary 12.2 --- Finite homological
palette}\label{corollary-12.2-finite-homological-palette}

By partition decomposition,

\[
H_q(K_r)
\simeq
\bigoplus_{k=1}^d
H_q(\operatorname{Conf}_k)^{\oplus S(r,k)}
\]

as abelian groups.

Therefore the \(E^1\)-page is built from the homology of only the
finitely many fiber configuration varieties

\[
\operatorname{Conf}_1,\dots,\operatorname{Conf}_d.
\]

The difficult information is not the list of groups appearing at higher
simplicial degrees; it is the system of face maps among their copies.

This gives a concrete program: \textbf{compute the finite configuration
palette and the finite set of induced maps on homology, then compute the
quotient topology.}

\begin{center}\rule{0.5\linewidth}{0.5pt}\end{center}

\section{13. Canonical commuting vector
fields}\label{canonical-commuting-vector-fields}

The étale structure gives more than a trivial tangent bundle.

Let \(y_1,\dots,y_n\) be the standard coordinates on the target. There
are canonical vector fields \(D_1,\dots,D_n\) on the source defined by

\[
D_i(f_j)=\delta_{ij}.
\]

Because the Jacobian determinant is a nonzero constant, the inverse
Jacobian has polynomial entries, so these are polynomial derivations.

They are the pullbacks of the coordinate vector fields on the target.
Hence:

\[
[D_i,D_j]=0,
\]

and at every point they form a basis of the tangent space.

The same construction works on every configuration component via any
étale projection to \(\mathbb A^n\).

\begin{center}\rule{0.5\linewidth}{0.5pt}\end{center}

\subsection{Proposition 13.1 --- Local-nilpotence criterion for
triviality}\label{proposition-13.1-local-nilpotence-criterion-for-triviality}

If all canonical derivations \(D_1,\dots,D_n\) of a connected Keller
\(n\)-fold are locally nilpotent, then the variety is isomorphic to
\(\mathbb A^n\).

For the original Keller map \(f\), if all canonical \(D_i\)'s are
locally nilpotent, then \(f\) is a polynomial automorphism.

\subsubsection{Proof}\label{proof-6}

Commuting locally nilpotent derivations integrate to an algebraic
\(\mathbb G_a^n\)-action. Since the vector fields form a basis
everywhere, every orbit has dimension \(n\), hence is open.
Connectedness forces a single orbit. The stabilizer is a
zero-dimensional algebraic subgroup of \(\mathbb G_a^n\), hence trivial
in characteristic zero. Thus the action is simply transitive and the
variety is \(\mathbb G_a^n\simeq\mathbb A^n\).

For the original \(f\), the identities \(D_i(f_j)=\delta_{ij}\) make
\(f\) translation-equivariant, so it is an isomorphism.

This gives a second bridge between Keller geometry and the
Makar--Limanov/LND viewpoint.

A counterexample must possess a global algebraic frame that is
\textbf{not complete in the additive algebraic sense}.

\begin{center}\rule{0.5\linewidth}{0.5pt}\end{center}

\section{14. Specialization: the marked-cubic Keller
map}\label{specialization-the-marked-cubic-keller-map}

Now specialize to the explicit degree-three map \(F\) already studied in
the current manuscript.

The established properties are:

\begin{enumerate}
\def\labelenumi{\arabic{enumi}.}
\tightlist
\item
  \(F\) has generic degree \(3\);
\item
  its fibers have cardinality exactly \(3,1,0\);
\item
  the nonproperness locus is the cubic discriminant hypersurface
  \(\mathcal S\);
\item
  the omitted locus is the triple-root curve \(\Gamma\);
\item
  the finite étale cover over \(\mathbb A^3\setminus\mathcal S\) has
  monodromy \(S_3\);
\item
  \(Y=\operatorname{Conf}_2(F)\) is a connected smooth affine threefold
  and a degree-six finite étale cover of the proper locus.
\end{enumerate}

\begin{center}\rule{0.5\linewidth}{0.5pt}\end{center}

\subsection{Theorem 14.1 --- Exceptional collision
collapse}\label{theorem-14.1-exceptional-collision-collapse}

For this \(F\),

\[
\operatorname{Conf}_1(F)=\mathbb A^3,
\qquad
\operatorname{Conf}_2(F)=Y,
\qquad
\operatorname{Conf}_3(F)\simeq Y,
\qquad
\operatorname{Conf}_k(F)=\varnothing\ (k\ge4).
\]

\subsubsection{Proof}\label{proof-7}

The only nontrivial point is
\(\operatorname{Conf}_3\simeq\operatorname{Conf}_2\).

The top forgetful map

\[
\operatorname{Conf}_3\to\operatorname{Conf}_2
\]

is an open immersion by Proposition 6.2. The fiber census has no
two-point fibers. Therefore every ordered pair of distinct points in one
fiber lies in a three-point fiber and has a unique third point. The open
immersion is surjective and hence an isomorphism.

Geometrically, this says that after two distinct roots of a cubic have
been chosen, the third is forced.

\begin{center}\rule{0.5\linewidth}{0.5pt}\end{center}

\subsection{\texorpdfstring{Corollary 14.2 --- Every level uses only
\(\mathbb A^3\) and
\(Y\)}{Corollary 14.2 --- Every level uses only \textbackslash mathbb A\^{}3 and Y}}\label{corollary-14.2-every-level-uses-only-mathbb-a3-and-y}

For every \(r\ge1\),

\[
K_r(F)
\simeq
\mathbb A^3
\sqcup
\left(S(r,2)+S(r,3)\right)Y.
\]

Using

\[
S(r,2)+S(r,3)=\frac{3^r-3}{6},
\]

we obtain the especially simple formula

\[
\boxed{
K_r(F)
\simeq
\mathbb A^3
\sqcup
\frac{3^r-3}{6}\,Y.
}
\]

Thus the entire infinite Keller nerve has only \textbf{two}
connected-component isomorphism types.

For every additive invariant \(I\),

\[
\boxed{
I(K_r(F))
=
I(\mathbb A^3)
+
\frac{3^r-3}{6}\,I(Y).
}
\]

Since the existing calculation gives

\[
\chi(Y)=0,
\]

it follows immediately that

\[
\boxed{
\chi(K_r(F))=1
\qquad\text{for every }r\ge1.
}
\]

\begin{center}\rule{0.5\linewidth}{0.5pt}\end{center}

\subsection{Proposition 14.3 --- Strong finite morphism type in the
degree-three
example}\label{proposition-14.3-strong-finite-morphism-type-in-the-degree-three-example}

The nontrivial simplicial structure maps can be expressed using a finite
collection of morphisms involving \(Y\):

\begin{itemize}
\tightlist
\item
  the \(S_3\) deck transformations of the full-ordering cover \[
  Y\to\mathbb A^3\setminus\mathcal S;
  \]
\item
  the three maps \[
  Y\to\mathbb A^3
  \] selecting the first, second, or uniquely determined third sheet;
\item
  the identity/diagonal maps on \(\mathbb A^3\).
\end{itemize}

Indeed, every component of every \(K_r\) is specified by a partition
with one, two, or three blocks. Deleting a coordinate either preserves
at least two distinct blocks, in which case the resulting map between
\(Y\)-components is a relabeling of the three sheets, or reduces to one
block, in which case it selects one sheet.

Thus this example is not merely finite-object-type; it is an unusually
small \textbf{finite-state simplicial system}.

This is the precise sense in which the symmetric marked-root geometry
causes the Keller groupoid to collapse.

\begin{center}\rule{0.5\linewidth}{0.5pt}\end{center}

\section{15. Why the symmetric construction first appears in dimension
three}\label{why-the-symmetric-construction-first-appears-in-dimension-three}

The projective construction publicly proposed by Andy Jiang on July 20,
2026 is the marked-root insertion map

\[
\mathbb P^1\times\operatorname{Sym}^2(\mathbb P^1)
\longrightarrow
\operatorname{Sym}^3(\mathbb P^1),
\qquad
(p,\{q,r\})\mapsto\{p,q,r\}.
\]

It is generically three-to-one because a generic cubic has three choices
of marked root.

More generally, the analogous operation

\[
\mathbb P^1\times\operatorname{Sym}^{m-1}(\mathbb P^1)
\to
\operatorname{Sym}^m(\mathbb P^1)
\]

is generically degree \(m\) and lives in dimension \(m\).

For \(m=2\), the analogous marked-root map would have generic degree
two. But a degree-two function-field extension in characteristic zero is
separable and normal, and Campbell's normal-extension theorem excludes a
noninvertible degree-two Keller map.

So the marked-root mechanism cannot produce a planar counterexample in
its naïve dimension-two form.

Dimension three is the first place where:

\begin{itemize}
\tightlist
\item
  the natural marked-root degree is at least three; and
\item
  the generic extension can be nonnormal;
\item
  \(S_3\) provides the smallest nonregular transitive monodromy.
\end{itemize}

This does \textbf{not} prove the planar Jacobian conjecture: a
hypothetical planar counterexample could have generic degree \(d\ge3\)
and arise from entirely different geometry.

But it explains why the symmetric construction itself begins naturally
in dimension three.

For provenance, the symmetric-product interpretation should be
attributed to Andy Jiang's July 20, 2026 public post, with the same-day
Secret Blogging Seminar thread used as contemporaneous corroboration.
Later papers should be cited for proof and exposition rather than
retroactively for priority; see \texttt{CITATION\_HYGIENE.md}.

\begin{center}\rule{0.5\linewidth}{0.5pt}\end{center}

\section{16. Imported global constraints on Keller
groupoids}\label{imported-global-constraints-on-keller-groupoids}

The following are useful existing results and should be treated as
literature inputs, not new theorems of this draft.

\subsection{16.1 Image topology}\label{image-topology}

Peretz--Van Chau--Campbell--Gutierrez record the classical facts that
for a complex Keller map

\[
f\colon\mathbb C^n\to\mathbb C^n,
\]

the image \(U=f(\mathbb C^n)\) is open and simply connected, while the
complement has complex codimension at least two.

In dimension two this means

\[
U=\mathbb C^2\setminus Z
\]

for a finite set \(Z\).

Reference:

R. Peretz, N. Van Chau, L. A. Campbell, C. Gutierrez, \emph{Iterated
images and the plane Jacobian conjecture}, Discrete Contin. Dyn. Syst.
16 (2006), 455--461, DOI \texttt{10.3934/dcds.2006.16.455}.

\begin{center}\rule{0.5\linewidth}{0.5pt}\end{center}

\subsection{16.2 Nonproperness
hypersurface}\label{nonproperness-hypersurface}

Jelonek proved that for a dominant generically finite polynomial map of
\(\mathbb C^n\), the nonproperness set is either empty or a hypersurface
(with strong additional uniruledness properties).

Reference:

Z. Jelonek, \emph{The set of points at which a polynomial map is not
proper}, Ann. Polon. Math. 58 (1993), 259--266, DOI
\texttt{10.4064/ap-58-3-259-266}.

Hence a nonautomorphic Keller map has two distinct exceptional objects:

\begin{itemize}
\tightlist
\item
  the omitted set \(Y\setminus U\), of codimension at least two;
\item
  the nonproperness hypersurface \[
  S_f=Y\setminus U_d.
  \]
\end{itemize}

The marked-cubic example displays this separation explicitly: the
omitted set is the curve \(\Gamma\), while \(S_f\) is the whole
discriminant hypersurface.

\begin{center}\rule{0.5\linewidth}{0.5pt}\end{center}

\subsection{16.3 Further planar
restrictions}\label{further-planar-restrictions}

Jelonek's 2020 note proves additional restrictions on the nonproperness
set of a hypothetical planar Keller counterexample, including that in
dimension two the nonproperness curve cannot be a curve without
self-intersections.

Reference:

Z. Jelonek, \emph{A note on the Jacobian Conjecture}, arXiv:2011.03472.

These results should be incorporated into the planar Keller-groupoid
program as constraints on the base \(V=\mathbb A^2\setminus S_f\) of the
finite étale monodromy cover.

\begin{center}\rule{0.5\linewidth}{0.5pt}\end{center}

\section{17. A Keller-groupoid formulation of the planar
problem}\label{a-keller-groupoid-formulation-of-the-planar-problem}

Suppose, hypothetically, that

\[
f\colon\mathbb A^2\to\mathbb A^2
\]

is a noninvertible Keller map of generic degree \(d\).

Then the Keller-groupoid framework forces the following package.

\subsection{17.1 Immediate structure}\label{immediate-structure}

\begin{enumerate}
\def\labelenumi{\arabic{enumi}.}
\tightlist
\item
  \(d\ge3\) (degree one is birational/invertible; degree two is excluded
  by the normal-extension criterion).
\item
  Every fiber configuration space \[
  \operatorname{Conf}_k(f),\quad 1\le k\le d,
  \] is a smooth affine algebraically parallelizable \textbf{surface}.
\item
  The full Čech nerve has only the finite palette \[
  \operatorname{Conf}_1,\dots,\operatorname{Conf}_d.
  \]
\item
  The proper locus \[
  V=\mathbb A^2\setminus S_f
  \] carries a finite étale degree-\(d\) cover.
\item
  The top configuration \(\operatorname{Conf}_d|_V\) is the Galois
  closure and its components are governed by the monodromy group \[
  G\le S_d.
  \]
\item
  The image \[
  U=f(\mathbb A^2)
  \] is the complement of finitely many points and is simply connected.
\end{enumerate}

This is already a highly constrained category of affine surfaces and
maps.

\begin{center}\rule{0.5\linewidth}{0.5pt}\end{center}

\section{18. Topological signal peculiar to dimension
two}\label{topological-signal-peculiar-to-dimension-two}

Let

\[
Z=\mathbb A^2\setminus U,
\qquad |Z|=\ell.
\]

Topologically,

\[
U^{an}=\mathbb C^2\setminus Z.
\]

For a finite set of \(\ell\) points, Alexander duality gives

\[
H_3(U;\mathbb Z)\simeq\mathbb Z^\ell,
\]

while

\[
H_1(U)=H_2(U)=0.
\]

On the other hand every level \(K_r(f)\) is a smooth affine complex
surface, so

\[
H_q(K_r)=0\qquad(q>2).
\]

Therefore, in the descent spectral sequence

\[
E^1_{p,q}=H_q(K_{p+1})
\Longrightarrow H_{p+q}(U),
\]

any nonzero \(H_3(U)\) must be manufactured \textbf{entirely by the
simplicial direction} from the rows \(q=0,1,2\).

Thus:

\begin{quote}
A hypothetical omitted point in the plane is visible as a genuinely
higher-simplicial homology class of the Keller groupoid; it cannot occur
inside the homology of any individual affine configuration surface.
\end{quote}

This does not yet yield a contradiction, but it identifies an explicit
topological target.

A sufficient new theorem of the following form would eliminate omitted
points:

\begin{quote}
\textbf{Desired planar vanishing theorem.}\\
For every planar Keller groupoid satisfying the algebraic constraints
above, the totalized collision complex has \(H_3=0\).
\end{quote}

If proved, this would force \(Z=\varnothing\), i.e.~surjectivity.

Surjectivity alone is not currently enough to conclude the planar
Jacobian conjecture, so a second step would still be required to
eliminate nontrivial surjective étale self-maps.

\begin{center}\rule{0.5\linewidth}{0.5pt}\end{center}

\section{19. Stable images and a sharper planar
reformulation}\label{stable-images-and-a-sharper-planar-reformulation}

The 2006 theorem of Peretz--Van Chau--Campbell--Gutierrez gives a
particularly natural interface with the groupoid viewpoint.

For a planar Keller map, the iterated images

\[
\mathbb A^2\supseteq f(\mathbb A^2)\supseteq f^2(\mathbb A^2)\supseteq\cdots
\]

eventually stabilize to a cofinite, open, simply connected set

\[
\Omega.
\]

The restricted map

\[
f|_\Omega\colon\Omega\to\Omega
\]

is a \textbf{surjective étale endomorphism}.

They show that the planar Jacobian conjecture is equivalent to the
assertion that every such stabilized map is injective.

In Keller-groupoid language:

\begin{quote}
\textbf{JC(2) is equivalent to saying that every stabilized planar
Keller groupoid on a cofinite simply connected open subset of
\(\mathbb A^2\) is the unit groupoid.}
\end{quote}

This reframing does not prove the conjecture, but it converts the
residual problem into a statement about the nonexistence of a nontrivial
étale equivalence-relation groupoid on a very simple open surface.

This seems like the correct point at which to bring in:

\begin{itemize}
\tightlist
\item
  classification of smooth affine/quasi-affine surfaces;
\item
  algebraic parallelizability of the fiber configuration spaces;
\item
  topology of complements of plane curves;
\item
  log Kodaira dimension;
\item
  monodromy and Galois-closure surfaces.
\end{itemize}

\begin{center}\rule{0.5\linewidth}{0.5pt}\end{center}

\section{20. Potential planar attack
routes}\label{potential-planar-attack-routes}

The following are research directions, not established results.

\subsection{Route A --- Classify parallelizable configuration
surfaces}\label{route-a-classify-parallelizable-configuration-surfaces}

Every \(\operatorname{Conf}_k(f)\) is a smooth affine surface with

\[
T_{\operatorname{Conf}_k}\simeq\mathcal O^{\oplus2}.
\]

Moreover it admits several related étale maps arising from coordinate
projections and the common target.

A useful classification theorem of the form

\begin{quote}
sufficiently constrained smooth affine parallelizable surfaces étale
over \(\mathbb A^2\) must be \(\mathbb A^2\) or have prescribed boundary
topology
\end{quote}

could immediately constrain the possible \(\operatorname{Conf}_k\).

The marked-cubic threefold demonstrates that trivial tangent bundle
alone is far from enough in dimension three; dimension two may
nevertheless be much more rigid.

\begin{center}\rule{0.5\linewidth}{0.5pt}\end{center}

\subsection{Route B --- Use the Galois-closure
surface}\label{route-b-use-the-galois-closure-surface}

The top configuration

\[
\operatorname{Conf}_d|_V
\]

is the Galois closure of the proper finite étale cover over the
complement of the nonproperness curve.

Thus a hypothetical planar counterexample produces a smooth affine
parallelizable surface carrying a finite étale Galois map

\[
Z\to\mathbb A^2\setminus S_f
\]

with nonregular transitive monodromy \(G\).

One can ask whether surface classification, log Kodaira dimension, or
boundary intersection theory permits such a \(Z\) to extend
simultaneously to the nonproper étale configuration surface required by
the Keller groupoid.

\begin{center}\rule{0.5\linewidth}{0.5pt}\end{center}

\subsection{Route C --- Homological descent
obstruction}\label{route-c-homological-descent-obstruction}

Compute the normalized total chain complex of

\[
H_q(K_r),\qquad q=0,1,2.
\]

The partition theorem means only the homology of
\(\operatorname{Conf}_1,\dots,\operatorname{Conf}_d\) ever occurs. If
geometric/morphism stabilization forces this total complex to have
vanishing degree-three homology, omitted points are impossible.

For an early-collapsing Keller groupoid like the marked-cubic example,
this computation should be exceptionally explicit and can serve as a
model calculation.

\begin{center}\rule{0.5\linewidth}{0.5pt}\end{center}

\subsection{Route D --- Euler moment
constraints}\label{route-d-euler-moment-constraints}

The equations

\[
\chi(\operatorname{Conf}_k)
=
\sum_{m=k}^d(m)_k\chi(U_m)
\]

turn collision Euler characteristics into factorial moments of the
fiber-cardinality distribution.

In dimension two the omitted set is finite, so

\[
|U_0|=\chi(U_0)
\]

is an integer directly recoverable from these moments.

One can search for geometric sign or divisibility restrictions on
\(\chi(\operatorname{Conf}_k)\) for parallelizable affine surfaces that
make the required moment system impossible.

\begin{center}\rule{0.5\linewidth}{0.5pt}\end{center}

\subsection{Route E --- Iteration
amplification}\label{route-e-iteration-amplification}

If \(f\) has generic degree \(d>1\), then the iterate \(f^r\) has
generic degree \(d^r\).

The planar stable-image theorem eventually places all iterates inside
the same cofinite dynamical environment.

This suggests a general meta-strategy:

\begin{quote}
Find a Keller-groupoid invariant that must grow with generic degree
under iteration but is uniformly bounded by the topology/algebraic
geometry of a fixed cofinite planar stable image.
\end{quote}

Any such invariant would force \(d=1\).

Possible candidates include:

\begin{itemize}
\tightlist
\item
  complexity of Galois-closure configuration surfaces;
\item
  boundary divisor complexity;
\item
  ranks of selected homology groups;
\item
  number or type of configuration components;
\item
  log-geometric invariants of compactifications.
\end{itemize}

At present this is only a strategy template.

\begin{center}\rule{0.5\linewidth}{0.5pt}\end{center}

\subsection{Route F --- Local inertia along the nonproperness
curve}\label{route-f-local-inertia-along-the-nonproperness-curve}

Over

\[
V=\mathbb A^2\setminus S_f,
\]

the finite cover has monodromy

\[
\pi_1(V)\to G\le S_d.
\]

All nontrivial monodromy is generated by topology around the plane curve
\(S_f\), because the full image \(U\) is simply connected.

Thus the groupoid translates the planar Jacobian conjecture into a
question about whether a plane-curve complement can support a transitive
monodromy representation whose associated collision/Galois-closure
surfaces extend étale across exactly the required fiber-drop strata.

Jelonek's restrictions on the topology/singularity of \(S_f\) should be
inserted at this point.

This looks especially compatible with braid-group and plane-curve
techniques.

\begin{center}\rule{0.5\linewidth}{0.5pt}\end{center}

\section{21. A general ``configuration complexity''
program}\label{a-general-configuration-complexity-program}

The degree \(d\) alone is crude. The Keller groupoid suggests finer
invariants.

Possible invariants include:

\subsection{21.1 Geometric generation
depth}\label{geometric-generation-depth}

\[
g(f)=
\min\{m:\text{all configuration component types appear by }C_m\}.
\]

Always \(g(f)\le d\).

For the known counterexample,

\[
g(F)=2<3.
\]

\begin{center}\rule{0.5\linewidth}{0.5pt}\end{center}

\subsection{21.2 Component orbit profile}\label{component-orbit-profile}

Let

\[
o_k(f)=\#\pi_0(\operatorname{Conf}_k(f)).
\]

Over the proper locus this is

\[
o_k(f)
=
\#\bigl(G\backslash\operatorname{Inj}([k],[d])\bigr).
\]

Thus the sequence \(o_k\) measures the transitivity profile of the
monodromy group.

\begin{center}\rule{0.5\linewidth}{0.5pt}\end{center}

\subsection{21.3 Additive configuration
profile}\label{additive-configuration-profile}

For an additive invariant \(I\), record

\[
(I(\operatorname{Conf}_1),\dots,I(\operatorname{Conf}_d)).
\]

This finite vector determines \(I(K_r)\) for every \(r\).

\begin{center}\rule{0.5\linewidth}{0.5pt}\end{center}

\subsection{21.4 Morphism complexity}\label{morphism-complexity}

Record the isomorphism classes of configuration components together with
the restrictions of all forgetful/permutation maps between them.

This is the appropriate invariant if one wants to reconstruct the
homotopy type of the quotient rather than merely additive invariants of
each level.

\begin{center}\rule{0.5\linewidth}{0.5pt}\end{center}

\subsection{21.5 Discovery-complexity
hypothesis}\label{discovery-complexity-hypothesis}

The marked-cubic counterexample has extraordinarily low configuration
complexity:

\[
\operatorname{Conf}_1=\mathbb A^3,\qquad
\operatorname{Conf}_2\simeq\operatorname{Conf}_3\simeq Y.
\]

Its entire simplicial geometry is generated by two varieties and a
finite collection of maps.

This motivates, but does not prove, the heuristic:

\begin{quote}
Keller maps arising from low configuration complexity may be
disproportionately discoverable by symbolic/AI-guided search.
\end{quote}

This should remain a stated research hypothesis rather than a historical
explanation unless additional examples support it.

\begin{center}\rule{0.5\linewidth}{0.5pt}\end{center}

\section{22. Fillings attached to configuration
covers}\label{fillings-attached-to-configuration-covers}

The existing threefold work contains one further piece that should now
be viewed as \textbf{extra structure attached to the groupoid}, rather
than part of the definition.

For the marked-cubic map:

\begin{itemize}
\tightlist
\item
  \(Y=\operatorname{Conf}_2(F)\) is the ordered-double configuration
  cover of the proper target locus;
\item
  the \(S_3\) monodromy has a sign character;
\item
  adjoining a square root of the cubic discriminant produces a natural
  two-sheeted cover;
\item
  filling this cover across the discriminant divisor gives the smooth
  affine threefold \(X\);
\item
  \(X\) is rational, factorial, rigid, simply connected, and
  noncontractible.
\end{itemize}

This suggests a broader question:

\begin{quote}
Given a Keller groupoid and a finite representation or character of its
monodromy group, when does the associated finite étale cover of the
proper locus admit a natural affine filling across the nonproperness
divisor?
\end{quote}

The Rees-rigid Keller threefold would then be the first example of a
\textbf{representation-theoretic Keller-groupoid filling}.

This is currently an open program, but it gives the existing
discriminant filling a natural place in the generalized story.

\begin{center}\rule{0.5\linewidth}{0.5pt}\end{center}

\section{23. Proposed paper
architecture}\label{proposed-paper-architecture}

A Keller-groupoid paper could now be organized as follows.

\subsection{Part I --- General theory}\label{part-i-general-theory}

\begin{enumerate}
\def\labelenumi{\arabic{enumi}.}
\tightlist
\item
  Keller maps and the image \(U\).
\item
  The étale relation groupoid.
\item
  Čech nerve and categorical \(2\)-coskeletality.
\item
  Collision varieties.
\item
  Partition/Stirling decomposition.
\item
  Generic degree and finite configuration palette.
\item
  Fiber-cardinality filtration and nonproperness.
\item
  Collision tower and top open immersion.
\item
  Monodromy, transitivity, and Galois closure.
\item
  Additive invariants and recurrence.
\item
  Topological realization and descent.
\item
  Canonical parallelism and vector fields.
\end{enumerate}

\subsection{Part II --- The first nontrivial
example}\label{part-ii-the-first-nontrivial-example}

\begin{enumerate}
\def\labelenumi{\arabic{enumi}.}
\setcounter{enumi}{12}
\tightlist
\item
  Provenance and the marked-root cubic construction.
\item
  Degree \(3\), \(3/1/0\) fiber census, \(S_3\) monodromy.
\item
  \(\operatorname{Conf}_2=\operatorname{Conf}_3=Y\).
\item
  Exact formula \[
  K_r=\mathbb A^3\sqcup\frac{3^r-3}{6}Y.
  \]
\item
  Finite-state simplicial maps.
\item
  The discriminant/sign cover and its affine filling.
\item
  Rees presentation, rigidity, and topology of the filling.
\end{enumerate}

\subsection{Part III --- Constraints and open
problems}\label{part-iii-constraints-and-open-problems}

\begin{enumerate}
\def\labelenumi{\arabic{enumi}.}
\setcounter{enumi}{19}
\tightlist
\item
  Collision complexity.
\item
  Stabilization questions.
\item
  Planar Keller groupoids and JC(2).
\item
  Classification of parallelizable configuration surfaces.
\item
  Iteration amplification.
\item
  Monodromy/inertia constraints.
\item
  Representation-theoretic fillings.
\end{enumerate}

This architecture preserves all of the existing threefold results while
making the general theory the conceptual center.

\begin{center}\rule{0.5\linewidth}{0.5pt}\end{center}

\section{24. Claims proved in this
draft}\label{claims-proved-in-this-draft}

The following are proposed as \textbf{LLM-proved} claims in the
temporary mini-Hardy registry:

\begin{itemize}
\tightlist
\item
  \texttt{KG-01}: Keller relation groupoid / Čech nerve construction.
\item
  \texttt{KG-02}: every nerve level and fiber configuration space is a
  smooth affine parallelizable \(n\)-fold.
\item
  \texttt{KG-03}: partition decomposition by fiber configuration spaces.
\item
  \texttt{KG-04}: fiber cardinality is bounded by generic degree;
  \(\operatorname{Conf}_k=\varnothing\) for \(k>d\).
\item
  \texttt{KG-05}: finite geometric palette; every component type appears
  by level \(d\).
\item
  \texttt{KG-06}: fiber-cardinality filtration via images of fiber
  configuration spaces.
\item
  \texttt{KG-07}: top forgetful map
  \(\operatorname{Conf}_d\to \operatorname{Conf}_{d-1}\) is an open
  immersion and is an isomorphism iff there are no \((d-1)\)-point
  fibers.
\item
  \texttt{KG-08}: monodromy-orbit description of connected configuration
  components.
\item
  \texttt{KG-09}: top configuration realizes the Galois closure.
\item
  \texttt{KG-10}: Stirling/Grothendieck/additive-invariant formulas and
  recurrence.
\item
  \texttt{KG-11}: Euler factorial-moment identities.
\item
  \texttt{KG-12}: canonical commuting parallelizing vector fields and
  the LND triviality criterion.
\item
  \texttt{KG-F01}: for the marked-cubic example,
  \(\operatorname{Conf}_1=\mathbb A^3,\ \operatorname{Conf}_2\simeq\operatorname{Conf}_3\simeq Y\).
\item
  \texttt{KG-F02}: for the marked-cubic example, \[
  K_r\simeq\mathbb A^3\sqcup((3^r-3)/6)Y.
  \]
\item
  \texttt{KG-F03}: the marked-cubic nerve has strong finite morphism
  type generated by the \(S_3\)-relabelings and three sheet projections.
\end{itemize}

The topological-descent theorem and the global image/nonproperness
results are standard/literature inputs rather than newly proved claims.

\begin{center}\rule{0.5\linewidth}{0.5pt}\end{center}

\section{25. Open claims / research
targets}\label{open-claims-research-targets}

\begin{itemize}
\tightlist
\item
  \texttt{KG-O01}: classify possible configuration complexity of Keller
  maps.
\item
  \texttt{KG-O02}: determine whether early collision stabilization
  forces additional algebraic restrictions on \(f\).
\item
  \texttt{KG-O03}: classify smooth affine parallelizable surfaces that
  can occur as planar Keller configuration components.
\item
  \texttt{KG-O04}: prove a homological-descent obstruction strong enough
  to constrain planar Keller groupoids.
\item
  \texttt{KG-O05}: find an invariant growing under iteration but bounded
  on a stable cofinite planar image.
\item
  \texttt{KG-O06}: characterize monodromy representations realizable by
  planar Keller groupoids.
\item
  \texttt{KG-O07}: develop representation-theoretic affine fillings of
  collision covers.
\item
  \texttt{KG-O08}: test whether low configuration complexity correlates
  with discoverability or with special algebraic normal forms.
\end{itemize}

\begin{center}\rule{0.5\linewidth}{0.5pt}\end{center}

\section{References / source hygiene}\label{references-source-hygiene}

The general scheme-theoretic arguments above use standard results
including:

\begin{itemize}
\tightlist
\item
  Stacks Project Tag \texttt{02GE}: diagonal of an unramified morphism
  is open.
\item
  Stacks Project Tag \texttt{024T}: sections of unramified separated
  morphisms are open and closed.
\end{itemize}

Global Keller-map inputs:

\begin{itemize}
\tightlist
\item
  L. A. Campbell, \emph{A condition for a polynomial map to be
  invertible}, Math. Ann. 205 (1973), 243--248.
\item
  Z. Jelonek, \emph{The set of points at which a polynomial map is not
  proper}, Ann. Polon. Math. 58 (1993), 259--266.
\item
  R. Peretz, N. Van Chau, L. A. Campbell, C. Gutierrez, \emph{Iterated
  images and the plane Jacobian conjecture}, DCDS 16 (2006), 455--461.
\item
  Z. Jelonek, \emph{A note on the Jacobian Conjecture},
  arXiv:2011.03472.
\end{itemize}

For the July 2026 marked-root example and priority-sensitive
attributions, use \texttt{CITATION\_HYGIENE.md}. In particular, the
symmetric-product marked-root interpretation should be attributed to
Andy Jiang's July 20 public post as the earliest public source currently
identified, with the Secret Blogging Seminar thread as contemporaneous
corroboration and later papers cited for detailed proofs/exposition.

\begin{center}\rule{0.5\linewidth}{0.5pt}\end{center}

\section{Immediate next calculations}\label{immediate-next-calculations}

\begin{enumerate}
\def\labelenumi{\arabic{enumi}.}
\tightlist
\item
  Write the explicit groupoid structure maps on
  \(Y=\operatorname{Conf}_2(F)\), including the third-sheet map
  \(Y\to\mathbb A^3\) and the full \(S_3\) deck action.
\item
  Compute \(H_*(Y)\), not only \(\chi(Y)\).
\item
  Build the actual \(E^1\) page of the Keller-groupoid descent spectral
  sequence for \(F\).
\item
  Verify directly that the finite-state description reproduces the
  topology of \[
  U=\mathbb A^3\setminus\Gamma.
  \]
\item
  Investigate smooth affine parallelizable surfaces as candidate planar
  fiber configuration varieties.
\item
  Recast the Peretz--Van Chau--Campbell--Gutierrez stable-image theorem
  entirely in Keller-groupoid language and look for an iteration-growth
  obstruction.
\end{enumerate}

\begin{center}\rule{0.5\linewidth}{0.5pt}\end{center}

\section{26. Plane-specific configuration
criteria}\label{plane-specific-configuration-criteria}

The general theory above becomes sharper in dimension two. The detailed
statements and proofs are collected in
\texttt{PLANE\_CONFIGURATION\_CRITERIA.md}.

The main new structural points are:

\begin{enumerate}
\def\labelenumi{\arabic{enumi}.}
\tightlist
\item
  \textbf{fiber deficit decomposition} \[
  d-m_C=\operatorname{supp}(\sigma_C)+t_C;
  \]
\item
  \textbf{deficit one is purely non-maximal} --- a \((d-1)\)-point curve
  stratum cannot come from branch ramification;
\item
  \textbf{étale-maximal top stabilization} \[
  \operatorname{Conf}_d(f)\simeq\operatorname{Conf}_{d-1}(f);
  \]
\item
  the configuration filtration recovers the generic deletion counts once
  inertia is known;
\item
  the complement divisors of configuration images give lower bounds on
  unit ranks of the configuration surfaces;
\item
  the first groupoid level admits an explicit divided-difference
  idempotent and ideal-membership criterion.
\end{enumerate}

These are the groupoid-derived constraints to combine next with the
one-place-at-infinity and boundary-graph machinery.

